\documentclass[a4paper,12pt]{article}
\usepackage[utf8]{inputenc}
\usepackage[italian]{babel}
\usepackage{listings}
\usepackage{xcolor}
\usepackage{float}
\usepackage{placeins}
\usepackage{url}
\setlength{\emergencystretch}{3em}

\lstset{
  language=C,
  backgroundcolor=\color{lightgray!20},
  frame=single,
  breaklines=true,
  numbers=left,
  numberstyle=\tiny,
  stepnumber=1,
  numbersep=5pt,
  captionpos=b,
  basicstyle=\ttfamily\small,
  keywordstyle=\color{blue}\bfseries,
  commentstyle=\color{green!60!black},
  stringstyle=\color{red},
  showstringspaces=false
}

\title{LabReSiD26 -- Hands-On 15}
\author{Davide Ferrara -- 518629}
\date{}

\begin{document}

\maketitle

% ---------------------------------------------------------------------------
\section{Esercizio 1: Buffer circolare a singolo thread}
% ---------------------------------------------------------------------------

\subsection*{Cosa succede se \texttt{buf\_get} viene chiamata su un buffer vuoto?}

Né \texttt{buf\_get} né \texttt{buf\_put} controllano \texttt{count} prima
di operare: il campo è un semplice contatore senza guardia. Chiamare
\texttt{buf\_get} su un buffer vuoto non produce nessun crash perché
\texttt{head} rimane sempre dentro i limiti dell'array grazie al modulo
$N$; \texttt{head} avanza circolarmente e torna a $0$, leggendo i valori
stantii lasciati dal round precedente. Nel frattempo \texttt{count} scende
a valori negativi senza alcun segnale di errore.

\begin{lstlisting}[language=bash, caption={Cinque \texttt{buf\_get} extra su buffer vuoto}]
Buffer: 12 13 14 15
Tutti gli elementi prelevati correttamente.
Buffer: 12 13 14 15
Elemento extra: 12
Elemento extra: 13
Elemento extra: 14
Elemento extra: 15
Elemento extra: 12
Buffer: 12 13 14 15
\end{lstlisting}

Le prime quattro chiamate restituiscono $12, 13, 14, 15$: i valori
dell'ultimo round ancora presenti nelle celle. La quinta restituisce $12$
di nuovo perché \texttt{head} ha completato il giro circolare e riparte
da $0$.

In un'implementazione corretta \texttt{buf\_get} dovrebbe controllare
\texttt{count} prima di leggere: se \texttt{count} $= 0$ tutti gli
elementi sono già stati prelevati e l'operazione non è consentita.
In questa versione \texttt{count} non viene mai controllato da nessuna
funzione, quindi il buffer non ha alcuna protezione contro letture oltre
il limite.

% ---------------------------------------------------------------------------
\section{Esercizio 2: Due thread senza sincronizzazione}
% ---------------------------------------------------------------------------

\begin{lstlisting}[language=bash, caption={Output Helgrind -- race condition in \texttt{buf\_put}}]
==21199== Possible data race during write of size 4 at 0x1FFEFFE4E8 by thread #2
==21199== Locks held: none
==21199==    at 0x40011D6: buf_put (in /home/dff/git/reti/ho15/step2)
==21199==    by 0x4001268: produttore (in /home/dff/git/reti/ho15/step2)
==21199==    by 0x485F47B: mythread_wrapper (hg_intercepts.c:410)
==21199==    by 0x49381B8: start_thread (pthread_create.c:454)
==21199==    by 0x49BD043: clone (clone.S:100)
\end{lstlisting}

% TODO

\end{document}
